\documentclass[11pt]{article}
\usepackage{amsmath,amssymb,amsthm}
\newtheorem{theorem}{Theorem}
\newcommand{\E}{\mathbb E}
\newcommand{\Cov}{\operatorname{Cov}}
\newcommand{\Var}{\operatorname{Var}}
\title{An expected ordinary covariance bound for the actual AI36 fiber}\author{}\date{}
\begin{document}\maketitle
Use the original isotropic Wilson measure, fixed $SU(N)$, and finite embedded neighborhoods.
Write $L_\beta=\log(2+\beta)$ and assume $\beta\ge1$. Constants depend only on the group
and fixed local geometry. This result is about the ACTUAL local conditional covariance,
not the variance of the full smeared observable. It supplies a quantitative bound for the
particular AI36 transition identified in the actual AI36 interface calculation.

\paragraph{Two inputs already proved for the same measure.}
The Haar total-force identities in the uniform contact moments calculation imply, uniformly for integers $p\ge1$,
\[
 \|\mathbf F_b\|_{2p}\le C\sqrt{p/\beta},\qquad \mathbf F_b=(X_b^AS_0)_A.
\]
The manuscript's elementary plaquette estimate PD12.31 gives
$\E s_p\le C L_\beta/\beta$. Its proof is reproduced in the orthogonal contact calculation and uses the original
partition function and isotropic symmetry. Therefore every individual plane-force vector
$A=\mathbf F_{b;\mu}$ is uniformly bounded and satisfies
\[
 |A|\le C,\qquad \E|A|^2\le C L_\beta/\beta.
\]
These statements alone do not bound its full fourth moment by $C\beta^{-2}$.

\paragraph{Exact resampling control of a directional force.}
Resample a single bond $e$ from its original one-link conditional law given all other links.
Call the original and resampled configurations $U,U'$. Both marginals are exactly $\mu$.
For every affected directional force $A=\mathbf F_{b;\mu}$,
\begin{equation}\label{delta}
 \|A(U)-A(U')\|_{2p}\le C\sqrt{p/\beta}.
\end{equation}
Here is the geometric reason, including the only exceptional row.
If $b\ne e$, two distinct bonds share at most one plaquette in the embedded neighborhood.
Thus either $A$ is unchanged, or its difference equals the difference of the TOTAL force
$\mathbf F_b$, since all other plaquette terms at $b$ are fixed.
If $b=e$, its plane force contains two plaquette terms. For either plaquette, choose its
opposite parallel bond $o\ne e$. The force vectors at $e$ and $o$ for this plaquette are
related by a minus sign and an adjoint transport along the perpendicular bath edge joining
their source vertices. This transport is independent of the resampled link.
For example, with plaquette word $P=U_e V U_o^{-1}W^{-1}$,
\[
 X_o s_p=-\operatorname{Ad}_{W^{-1}}X_e s_p
\]
when both gradients are identified with Lie-algebra vectors in AI1 normalization.
This follows by differentiating under the trace and cycling the inserted generator; the
sign changes with the orientation convention but its norm and independence do not.
The other incident plane plaquette has the same property with its own bath transport.
Each opposite bond shares just this plaquette with $e$, so its varying partial force equals
its varying total force. The central directional-force difference is consequently the sum
of at most two transported total-force differences at opposite bonds. Applying the total-force
moment bound to both marginals proves \eqref{delta}. No conditional concentration or
pointwise conditional gap has been assumed.

\paragraph{Square-force fluctuation lemma.}
For such a vector $A$, write $A'=A(U')$. Since
$||A|^2-|A'|^2|\le|A-A'|\,|A+A'|$, Holder with
$q=p/(p-1)$, $p\ge2$, gives
\begin{align*}
 \E\bigl(|A|^2-|A'|^2\bigr)^2
 &\le\|A-A'\|_{2p}^2\|A+A'\|_{2q}^2\\
 &\le C\frac p\beta\left(\frac{L_\beta}{\beta}\right)^{1-1/p}.
\end{align*}
The second inequality uses the uniform bound on $A+A'$ and
$\E|A+A'|^2\le C L_\beta/\beta$, with both original marginals retained.
Choose $p=\max(2,\lceil L_\beta\rceil)$. Since
$(\beta/L_\beta)^{1/p}\le C$, this yields the explicit bound
\begin{equation}\label{square}
 \E\bigl(|A|^2-|A'|^2\bigr)^2\le C\beta^{-2}L_\beta^2.
\end{equation}
For a total-force square, the sharper $C\beta^{-2}$ bound follows directly from its
fourth moment. The logarithmic loss in \eqref{square} records precisely where the
available directional-force energy estimate is used; no summable covariance assumption
has been inserted.

\begin{theorem}[Expected AI36 ordinary covariance]
Let $\Lambda$ consist of a separated packing of the two-link blocks of the orthogonal contact calculation, with
all their determining neighborhoods disjoint. Let $q$ observe every outside link and let
$\nu=q_*\mu$ be its actual marginal. Let $\mathcal R_j$ be the finite set of force rows
whose plaquette stars meet the free links in block $j$, including every original saved
pointing offset. For a literal test define
$W_j(f)^2=\sum_{b\in\mathcal R_j}|f_b|^2$ and $W(f)^2=\sum_jW_j(f)^2$.
For the original physical $G=O_H$ or the specified ordinary $G=r_H$,
\[
 \E_\nu\Var(G(f)\mid q)\le Cc_\beta^{-2}a^{-4}L_\beta^2 W(f)^2,
\]
and
\[
 \|\Cov(G(f),G(h)\mid q)\|_{L^1(\nu)}
 \le Cc_\beta^{-2}a^{-4}L_\beta^2 W(f)W(h).
\]
The same inequalities hold for the residual $S=G-R$ of the orthogonal contact calculation through the actual AI36 interface calculation.
For pointed completions replace $W_j(f)^2$ by the fixed finite sum of the actual
transformed test weights at their original rows. No common-center test replacement is made.
\end{theorem}
\begin{proof}
For one resampled link, only a fixed number of rows varies in $D[f]$ or $P[f]$.
Cauchy--Schwarz over those actual rows and \eqref{square} bound the expected squared
difference of either dimensionless sum by
$C\beta^{-2}L_\beta^2\sum_{b\text{ affected}}|f_b|^2$.
The physical factor is exactly $\sigma^2=g^{-4}a^{-4}$, so
$\sigma^2\beta^{-2}=c_\beta^{-2}a^{-4}$.
For $r_H$, the additional term is $-\sigma\ell[f]/(3\beta)$.
Its exact Casimir formula is a finite weighted plaquette sum. Compactness bounds its
resampling squared difference by $C\sigma^2\beta^{-2}$ times the sum of squared
original incident test weights. This is already at the claimed scale, with its constant
and all four incidences intact.

Conditioned on $q$, every free-link law factors. Tensorization of conditional variance
therefore gives
\[
 \Var(G\mid q)\le\sum_{e\in\Lambda}\E[\Var(G\mid U_{\ne e})\mid q].
\]
Each one-link variance is one half the expected squared difference under the conditional
independent replica. Average and insert the just-proved estimates. The number of rows,
colors, and plaquette incidences per link is fixed, and the packing keeps total overlap
bounded. This proves the first bound. Conditional Cauchy--Schwarz followed by ordinary
Cauchy--Schwarz gives the mixed $L^1$ bound.
Finally the orthogonal contact calculation and the local source transition calculation prove the stronger conditional orthogonal identity
\[
 \Var(G(f)\mid q)=v_{ff}(q)+\Var(S(f)\mid q),\qquad v_{ff}\ge0.
\]
Thus the same upper bound applies to the residual conditional variance and, by the same
Cauchy--Schwarz argument, to its mixed covariance. No independence of the outside bath
variables is required anywhere.
\end{proof}

\paragraph{What this advances and what it leaves.}
This is a linear support-count estimate, with two explicit logarithms, for the complete
ordinary covariance produced by the identified original compact conditional fiber.
For bounded physical tests $W(f)^2\le C\|f\|_\infty^2N_f$, with $N_f$ the literal
completed support plaquette count. The extracted contact itself has the sharper
$Ca^{-4}N_f$ bound without logarithms from the uniform contact moments calculation.
The law of total variance still contains
\[
 \Var_\mu G(f)=\E_\nu\Var(G(f)\mid q)+\Var_\nu\E(G(f)\mid q).
\]
The second term is the inherited boundary ordinary covariance and is not bounded here.
After an integration the next effective score is generally not the original Wilson score;
repeating the original force Ward recurrence without that score would be invalid.
Thus this is a reusable verified local transition estimate, not an all-history MT36 bound.
Bath-measurable selectors can be retained in unnormalized positive moment integrals;
normalizing a rare branch and free-link-dependent selectors need their own denominators
and complete responses. Complex domains, boundary-current norms, and the independent
absolute scalar factors are not supplied by this positive-measure bound.
\end{document}
